\documentclass{UnBExam}%

\disciplina{Estruturas de Dados}%
\turma{02}%
\documento{}%
\subdoc{Plano de Ensino}%
\periodo{2022/2}%
\professor{Guilherme N. Ramos}%116319 &

% ---------- Pacotes ---------- %
%\usepackage{listings}%             Código
%

% ---------- Estilo ---------- %
%\pagestyle{empty}%         remover o  número de página
% % \setlength\parindent{0pt}%  remover tabulação

% % % Mudar estilo da seção
% % \makeatletter%
% % \renewcommand\section{\@startsection{section}{1}{\z@}%
% %     {-3.5ex \@plus -1ex \@minus -.2ex}%
% %     {2.3ex \@plus.2ex}%
% %     {\normalfont\normalsize\bfseries}}%
% % \makeatother%

% % Novo estilo de enumeração
% \renewcommand{\theenumi}{\small\thesection.\arabic{enumi}}%
% \newenvironment{my_enumerate}%
% {\begin{enumerate}%
%   \vspace{-0.3\baselineskip}
%   \setlength{\itemsep}{1pt}%
%   \setlength{\parskip}{0pt}%
%   \setlength{\parsep}{0pt}}%
% {\end{enumerate}}%
% % Novo estilo para itemização
% \renewcommand{\labelitemi}{-}
% \newenvironment{my_itemize}%
% {\begin{itemize}%
%   \vspace{-0.3\baselineskip}
%   \setlength{\itemsep}{1pt}%
%   \setlength{\parskip}{0pt}%
%   \setlength{\parsep}{0pt}}%
% {\end{itemize}}

\printanswers
\begin{document}
    \section{Objetivos}%
      \vspace{-0.5\baselineskip}%
      Propiciar aos alunos conceitos sólidos na forma de armazenamento e organização dos dados por meio do gerenciamento dinâmico da memória e de estruturas de dados lineares e não lineares. Com a compreensão desses conceitos é esperado o desenvolvimento do pensamento computacional e das habilidades de programação com atividades práticas.

    \section{Programa}%
        \vspace{-0.5\baselineskip}%
        Gerenciamento dinâmico da memória. Listas. Pilhas. Filas. Árvores. Ordenação e Pesquisa.Pesquisa em Árvores. Grafos. Análise e aplicação de estruturas avançadas em problemas de programação. Resolução de problemas. Aplicações em problemas ambientais.

    \section{Metodologia}%
        \vspace{-0.5\baselineskip}%
        Leitura de material de referência para compreensão do conteúdo teórico.
        Aulas expositivas com parte do conteúdo teórico para reforço e esclarecimento.
        Aulas práticas e resolução de exercícios para estudo e fixação do conteúdo.
        Atendimento para os estudantes realizado pelos monitores, tutores e professores.

    \section{Avaliação da Aprendizagem}%
        \vspace{-0.5\baselineskip}%
        A frequência será registrada durante as aulas presenciais no decorrer da disciplina.
        A avaliação da aprendizagem será feita por meio de 2 provas presenciais (P) e realização de 1 projeto de programação (Prj). A nota final (NF) é definida por:
        $$MP = \frac{2P_1+3P_2}{5}$$
        $$NF = 0.6MP+0.4Prj$$

        A menção final é baseada na nota, conforme o regulamento da UnB. A aprovação está condicionada à uma nota igual ou superior a 50\% para cada componente da nota final (tanto MP quanto Prj). O aluno que não obtiver frequência mínima de 75\% em relação ao número total de aulas estará reprovado por faltas e receberá Menção Final SR independentemente do valor da Nota Final.

        \textbf{O aluno que copiar ou plagiar as atividades práticas ou falsificar a identificação em qualquer atividade será automaticamente reprovado na disciplina por questões éticas.}
        % Trabalhos $\left(M_T=\frac{2T_1+3T_2}{5}\right)$
    %   \vspace{-0.5\baselineskip}%
    %   \begin{my_enumerate}
    %       \item Do Aluno:\vspace{-.6\baselineskip}\\
    %           \begin{minipage}{.5\textwidth}
    %           \begin{my_itemize}
    %               \item 2 Trabalhos $\left(M_T=\frac{2T_1+3T_2}{5}\right)$
    %           \end{my_itemize}
    %           \end{minipage}
    %           \begin{minipage}{.5\textwidth}
    %           \begin{my_itemize}
    %               \item 3 Provas $\left(M_P= \overbrace{\frac{1.1P_{\alpha} + P_{\beta} + 0.9P_{\gamma}}{3}}^{P_{\alpha} \leq P_{\beta}\leq P_{\gamma}}\right)$
    %           \end{my_itemize}
    %           \end{minipage}
    %           \begin{minipage}{.15\textwidth}
    %           \end{minipage}
    %           \vspace{-1em}
    %           \begin{my_itemize}
    %               \item Média: $0.75M_P+0.25M_T$
    %           \end{my_itemize}
    %       \item Da Disciplina:
    %           \begin{my_itemize}
    %               \item Questionário de avaliação da UnB ao final do semestre.
    %           \end{my_itemize}
    %   \end{my_enumerate}

    \section{Bibliografia}%
        \vspace{-0.5\baselineskip}%
        \begin{enumerate}
            \item Brad Miller e David Ranum. Resolução de Problemas com Algoritmos e Estruturas de Dados usando Python, 2013.
            \vspace{-0.5\baselineskip}%
            \item Allen B. Downey. Pense em Python, 2a Edição, 2016.
            \vspace{-0.5\baselineskip}%
            \item Brad Miller e David Ranum. Como Pensar Como um Cientista da Computação, 2008.
        \end{enumerate}
    %   \vspace{-0.5\baselineskip}
    %   \begin{my_enumerate}
    %       \item Nivio Ziviani. \emph{Projeto de Algoritmos com implementações em Pascal e C}. Pioneira Thomson Learning, 2005
    %       \item Aaron Tenenbaum, Yedidyah Langsam, Moshe J. Augenstein. \emph{Estruturas de Dados usando C}. Pearson Makron Books, 1995
    %       \item Cormen, Leiserson, Rivest, Stein \emph{Introduction to Algorithms}. The MIT Press, 2009
    %   \end{my_enumerate}

    % \section{Observações}%
    %   \vspace{-0.5\baselineskip}
    %   \begin{my_enumerate}
    %       \item \textbf{As médias devem ser iguais ou superiores a 5 ($M_T \geq5$ e $M_P \geq5$).}
    %       \item Calendário de provas: 3 de maio, 7 de junho e 19 de julho.
    %       \item Não serão realizadas provas de reposição, em nenhuma hipótese. Casos específicos serão analisados pelo Professor.
    %   \end{my_enumerate}

    \vspace{2em}%\vfill
    Brasília, \today \hfill \parbox[t]{75mm} {\hrulefill\\\ \ Prof.}
\end{document}